\section{Discriminant of a sesquilinear form}

Throughout this paragraph, A denotes a commutative ring, J an automorphism of A, and E a free A-module of finite dimension n.

DEFINITION 1. --- Given a form \(\Phi\) sesquilinear for J on E and a system \(S = (x_1, \ldots, x_n)\) of n elements of E, one calls the discriminant of \(\Phi\) with respect to this system, and denotes it by \(D_{\Phi}(x_1, \ldots, x_n)\) or \(D_{\Phi}(S)\) the element det \((\Phi(x_i, x_j))\) of A.

If \((e_{1},\ldots,e_{n})\) is a basis of E, the discriminant of \(\Phi\) with respect to this basis is none other than the determinant of the matrix of \(\Phi\) with respect to this basis.

It follows from the definition of the extension of \(\Phi\) to \(\wedge\) E ( \(\S\) 1, no. 9) that we have

\[
\mathrm{D} _ {\Phi} (x _ {1}, \dots , x _ {n}) = \Phi_ {(n)} (x _ {1} \wedge \dots \wedge x _ {n}, x _ {1} \wedge \dots \wedge x _ {n}),\tag{1}
\]

where \(\Phi_{(n)}\) denotes the extension of \(\Phi\) to \(\wedge\) E. For every permutation \(\sigma \in \mathfrak{S}_n\) , we therefore have

\[
\mathrm{D} _ {\Phi} (x _ {\sigma (\mathbf {1})}, \dots , x _ {\sigma (n)}) = \mathrm{D} _ {\Phi} (x _ {\mathbf {1}}, x _ {\mathbf {2}}, \dots , x _ {n}).
\]

Example. --- Let B be an algebra over the ring A, such that B is a free A-module of finite dimension n. Then the map \((x, y) \to \mathrm{Tr}_{\mathrm{B}/\mathrm{A}}(xy)\) (chap.~VIII, § 12, no. 2) is a bilinear form on B. Given a system \((x_{1}, \ldots, x_{n})\) of n elements of B, the discriminant of this form with respect to this system is called the discriminant of the system \((x_{1}, \ldots, x_{n})\) over A, and is denoted \(\mathrm{D}_{\mathrm{B}/\mathrm{A}}(x_{1}, \ldots, x_{n})\) . We thus have

\[
\mathrm{D} _ {\mathrm{B/A}} (x _ {1}, \dots , x _ {n}) = \det (\operatorname{Tr} _ {\mathrm{B/A}} (x _ {i} x _ {j})).\tag{2}
\]

Remark. --- Let \((e_{1},\ldots,e_{n})\) be a basis of B over A, and \(e_{i}e_{j}=\sum_{k=1}^{n}c_{ijk}e_{k}\;(c_{ijk}\in\mathrm{A})\) the multiplication table of this basis (chap.~II, § 7, no. 2). Since the matrix of the A-linear endomorphism \(x\to e_{k}x\) of B with respect to \((e_{r})\) is \((c_{ksr})\) , one has \(\mathrm{Tr}_{\mathrm{B}/\mathrm{A}}(e_{k})=\sum_{r=1}^{n}c_{krr}\) , whence \(\mathrm{Tr}_{\mathrm{B}/\mathrm{A}}(e_{i}e_{j})=\sum_{k,r}c_{ijk}c_{krr}\) . It follows that one has

\[
\mathrm{D} _ {\mathbf {B} / \mathbf {A}} (e _ {1}, \dots , e _ {n}) = \det _ {i, j} (\sum_ {k, r} c _ {i j k} c _ {k r r}).\tag{3}
\]

PROPOSITION 1. — Let \(\Phi\) be a sesquilinear form for J on E, \((x_1,\ldots,x_n)\) be a system of \(n\) elements of E, and \((a_{ij})_{(i,j=1,\ldots,n)}\) a family of \(n^2\) elements of A; set \(y_j = \sum_{i=1}^{n} a_{ji} x_i\) . We then have

\[
\mathrm{D} _ {\Phi} (y _ {1}, \dots , y _ {n}) = \det (a _ {i j}). \det (a _ {i j}) ^ {\mathrm{J}}. \mathrm{D} _ {\Phi} (x _ {1}, \dots , x _ {n}).
\]

Indeed, as \(\Phi\) is a sesquilinear form, we have \(\Phi(y_i, y_j) = \sum_{k,m} a_{ik} \Phi(x_k, x_m) a_{jm}^{\mathfrak{J}}\) . Therefore, if we denote by \(A\) the matrix \((a_{ij})\) , the matrix \((\Phi(y_i, y_j))\) is equal to \(A \cdot (\Phi(x_i, x_j)) \cdot {}^t A^{\mathfrak{J}}\) . Since \(\det(^t A) = \det(A)\) and \(\det(A^{\mathfrak{J}}) = \det(A)^{\mathfrak{J}}\) , our assertion is proved.

In particular, if \((e_{i})\) and \((e_{i}^{\prime})\) are two bases of E, D and D' the discriminants of \(\Phi\) with respect to these bases, and a the determinant of the change-of-basis matrix from the basis \((e_{i})\) to the basis \((e_{i}^{\prime})\) , one has

\[
\mathrm{D} ^ {\prime} = a a ^ {\mathrm{J}} \mathrm{D}.\tag{4}
\]

It follows from prop. 1 that, if \((e_i)\) is a basis of E and \((x_i)\) any arbitrary system of \(n\) elements of E, \(\mathrm{D}_{\Phi}(e_1,\ldots ,e_n)\) divides \(\mathrm{D}_{\Phi}(x_1,\dots,x_n)\) . In particular, the discriminants of \(\Phi\) with respect to any two bases of E generate the same principal ideal of A.

Let \((\mathbf{E}_i)_{i\in I}\) be a finite family of free A-modules of finite dimensions, \(\Phi_i\) a sesquilinear form for J on \(\mathbf{E}_i\) , and \(\mathbf{B}_i\) a basis of \(\mathbf{E}_i\) . If \(\Phi\) denotes the direct sum of \(\Phi_i\) ( \(\S 1\) , no. 3) and B the basis of \(\prod E_i\) obtained by taking the union of \(\mathbf{B}_i\) , we obviously have

\[
\mathrm{D} _ {\Phi} (\mathrm{B}) = \prod_ {i \in \mathrm{I}} \mathrm{D} _ {\Phi_ {i}} (\mathrm{B} _ {i}).\tag{5}
\]

Let \(\Phi\) be a sesquilinear form for \(J\) on \(E\) , \(h\) a homomorphism of \(A\) into a commutative ring \(A'\) , \(\Phi'\) the sesquilinear form on \(A' \otimes_A E\) obtained by extension of \(\Phi\) ( \(\S 1\) , no. 4) and \((x_1, \ldots, x_n)\) an arbitrary system of elements of \(E\) . Since \(A' \otimes_A E\) is an \(A'\) -free module, \(D_{\Phi'}(1 \otimes x_1, \ldots, 1 \otimes x_n)\) is defined, and we obviously have

\[
\mathrm{D} _ {\Phi^ {\prime}} (1 \otimes x _ {1}, \dots , 1 \otimes x _ {n}) = h (\mathrm{D} _ {\Phi} (x _ {1}, \dots , x _ {n})).\tag{6}
\]

Example. --- Let B be an algebra over A which is a free A-module of finite dimension \(n\) , \((x_{1},\ldots,x_{n})\) a basis of B over A, and m an ideal of A. If we denote by \(h\) the canonical homomorphism from B onto B/mB, \((h(x_{1}),\ldots,h(x_{n}))\) is a basis of B/mB over A/m (chap.~I, § 6, no. 5, prop. 5), and B/mB is isomorphic to \((A/m)\otimes_{A}B\) . We therefore have

\[
\mathrm{D} _ {(\mathrm{B} / \mathfrak {m B}) / (\mathrm{A} / \mathfrak {m})} (h (x _ {1}), \dots , h (x _ {n})) = h (\mathrm{D} _ {\mathrm{B} / \mathrm{A}} (x _ {1}, \dots , x _ {n})).
\]

PROPOSITION 2. --- Assume that A is an integral domain. Let \(\Phi\) be a sesquilinear form for J on E and \((e_1, \ldots, e_n)\) a basis of E, such that \(\mathrm{D}_{\Phi}(e_1, e_2, \ldots, e_n) \neq 0\) .

\begin{enumerate}
\def\labelenumi{\alph{enumi})}
\item
  For a system \((x_{1},\ldots ,x_{n})\) of \(n\) elements of E to be free, it is necessary and sufficient that \(\mathrm{D}_{\Phi}(x_1,\dots,x_n)\) be \(\neq 0\) .
\item
  For a system \((x_{1},\ldots ,x_{n})\) of \(n\) elements of E to be a basis of E, it is necessary and sufficient that \(\mathrm{D}_{\Phi}(x_1,\dots,x_n)\) and \(\mathrm{D}_{\Phi}(e_1,\dots,e_n)\) be associated elements in A (cf.~Chap.~VI, § 1, no.~5).
\end{enumerate}

Let us set \(x_j = \sum_{i=1}^{n} a_{ji} e_i\) ( \(a_{ji} \in A\) ). Let us first prove \(a\) ). If \(D_{\Phi}(x_1, \ldots, x_n) = 0\) , one has \(\det(a_{ji})\) . \(\det(a_{ji})^{\mathfrak{J}} = 0\) (prop. 1) since \(D_{\Phi}(e_1, \ldots, e_n) \neq 0\) and A is an integral domain; we therefore have \(\det(a_{ji}) = 0\) , and the vectors \(x_j\) are linearly dependent (chap.~III, § 7, no. 1, th. 1, applied to the vector space \(K \otimes_A E\) , where K denotes the field of fractions of A). Conversely, if these vectors are linearly dependent, we have \(\det(a_{ji}) = 0\) ( \(ibid.\) ), whence \(D_{\Phi}(x_1, \ldots, x_n) = 0\) (prop. 1).

We now prove \(b\) ). If \(\mathrm{D}_{\Phi}(x_1, \ldots, x_n)\) and \(\mathrm{D}_{\Phi}(e_1, \ldots, e_n)\) are associated in A, prop. 1 shows that \(\det(a_{ij}) \cdot \det(a_{ij})^{\mathrm{J}}\) is invertible in A. Thus \(\det(a_{ij})\) is also invertible in A; therefore the matrix \((a_{ij})\) on A is invertible (chap.~III, § 6, no. 5, th. 2), and the endomorphism \(g\) of E defined by \(g(e_i) = x_i\) ( \(i = 1, \ldots, n\) ) is an automorphism; consequently \((x_1, \ldots, x_n)\) is a basis of E. The converse follows immediately from Prop. 1.

PROPOSITION 3. --- Let \(\Phi\) be a sesquilinear form for J on E, and S a basis of E. The following conditions are equivalent:

\begin{enumerate}
\def\labelenumi{\alph{enumi})}
\item
  The map \(s_{\Phi}\) from E into \(E^{*}\) associated with \(\Phi\) is bijective.
\item
  The map \(d_{\Phi}\) from E into \(E^{*}\) associated with \(\Phi\) is bijective.
\item
  The element \(D_{\Phi}(S)\) is invertible in A.
\end{enumerate}

Indeed the condition \(c\) expresses that the matrix of \(\Phi\) with respect to S is invertible (chap. III, § 6, no. 5, th. 2). Therefore \(c\) is equivalent to \(a\) ( \(§ 1\) , no. 10); likewise \(c\) is equivalent to \(b\) ).

PROPOSITION 4. --- Assume A is an integral domain. Let S be a basis of E. A necessary and sufficient condition for a sesquilinear form \(\Phi\) on E to be nondegenerate is that one has D \(_{\Phi}\) (S) ≠ 0.

Indeed, let K be the field of fractions of A, and let \(\Phi'\) be the extension of \(\Phi\) to the K-vector space \(K\otimes_{\mathbf{A}}E\) ; we identify E with a subset of this vector space. The relation \(D_{\Phi}(S)\neq0\) is then equivalent to \(D_{\Phi'}(S)\neq0\) by formula (6), which itself expresses that \(s_{\Phi'}\) is bijective (prop. 3), that is, \(\Phi'\) is nondegenerate ( \(\S1\) , no. 6, prop. 6). Now, for every \(x\in K\otimes_{\mathbf{A}}E\) , there exists \(a\in A\) such that \(ax\in E\) . Therefore, for \(\Phi\) to be degenerate, it is necessary and sufficient that \(\Phi'\) be so. This proves our assertion.

PROPOSITION 5. --- Let A be a field, B a commutative algebra of finite dimension n over A, and S a basis of B. For B to be separable (chap. VIII, § 7, no. 5, def. 1) it is necessary and sufficient that one have D \(_{B/A}\) (S) ≠ 0.

Let A' be the algebraic closure of A, and let B' be the algebra A' ⊗\_A B over A'. If B is separable, B' is semisimple (Chap. VIII, § 7, no. 5, cor. of Prop. 7) and is therefore a direct product of n fields isomorphic to A' (Chap. VIII, § 6, no. 4, cor. of Prop. 9). If S' denotes the canonical basis of B' (identified with A'	extsuperscript{n}), we have D\_\{B'/A'\}(S') = 1, hence D\_\{B'/A'\}(S) ≠ 0 (prop. 1) and D\_\{B/A\}(S) ≠ 0 (formula (6)).

Conversely, suppose we have \(D_{B/A}(S) \neq 0\) . To show that B is separable, it suffices to show that \(B'\) is semisimple, that is, that it admits no nilpotent element \(\neq 0\) . Now, if \(x'\) were a nonzero nilpotent element of \(B'\) , one could take it as the first element of a basis \(S'\) of \(B'\) , and one would then have \(\mathrm{Tr}_{\mathbf{B}'/\mathbf{A}'}(x'y') = 0\) for all \(y' \in S'\) since a nilpotent endomorphism has all its eigenvalues zero (chap.~VII, § 5, no. 3, cor. 3 of prop. 8), hence zero trace. It would then follow that \(\mathrm{D}_{\mathbf{B}'/\mathbf{A}'}(\mathrm{S}') = 0\) , whence \(\mathrm{D}_{\mathbf{B}'/\mathbf{A}'}(\mathrm{S}) = 0\) (prop. 1) and \(\mathrm{D}_{\mathbf{B}/\mathbf{A}}(\mathrm{S}) = 0\) (formula (6)), contrary to the hypothesis.

Remark. --- Suppose that B is an overfield of A. Let \(S = (x_{1}, \ldots, x_{n})\) be a basis of B, and \((s_{1}, \ldots, s_{n})\) the A-isomorphisms of B into the algebraic closure A' of A (each \(s_{j}\) being repeated \([B : A]_{i}\) times). Recall that, for every \(z \in B\) , one has \(\operatorname{Tr}_{\mathbf{B}/\mathbf{A}}(z) = \sum_{j=1}^{n} s_{j}(z)\) (chap.~VIII, § 12, no. 2, prop. 4).

It then follows from the product formula for determinants that one has

\[
(\det (s _ {j} (x _ {i}))) ^ {2} = \mathrm{D} _ {\mathrm{B/A}} (x _ {1}, \dots , x _ {n}).\tag{7}
\]

This formula shows that proposition 5 generalizes the separability condition given in chap.~V, § 7, no. 2, Remark.

PROPOSITION 6. --- Let \(\Phi\) be an A-bilinear form on E, and K a subring of A such that A is a free K-module of finite dimension q. If \((e_i)_{i=1,\ldots,n}\) is a basis of E over A and \((a_j)_{j=1,\ldots,q}\) a basis of A over K, then \((a_j e_i)\) is a basis of E over K. The map \(\Phi'\) from E \(\times\) E to K defined by \(\Phi'(x,y)=\mathrm{Tr}_{\mathrm{A/K}}(\Phi(x,y))\) is a K-bilinear form on E, and one has

\[
\mathrm{D} _ {\Phi^ {\prime}} (a _ {j} e _ {i}) = \mathrm{N} _ {\mathrm{A} / \mathrm{K}} (\mathrm{D} _ {\Phi} (e _ {1}, \dots , e _ {n})) \cdot (\mathrm{D} _ {\mathrm{A} / \mathrm{K}} (a _ {1}, \dots , a _ {q})) ^ {n}.\tag{8}
\]

The first two assertions being evident, it suffices to prove (7). By definition, the left-hand side is the determinant of the K-linear endomorphism u of E defined by

\[
u (a _ {j} e _ {i}) = \sum_ {r, s} \mathrm{Tr} _ {\mathbf {A} / \mathbf {K}} (a _ {j} a _ {r} \Phi (e _ {i}, e _ {s})) a _ {r} e _ {s}.
\]

Consider the A-linear endomorphism \(\nu\) of E defined by \(\nu(e_i) = \sum_s \Phi(e_i, e_s)e_s\) , and the K-linear endomorphism \(w\) of E defined by \(w(a_j e_i) = (\sum_r \mathrm{Tr}_{\mathrm{A/K}}(a_j a_r)a_r)e_i\) . We have

\[
w \left(v \left(a _ {j} e _ {i}\right)\right) = w \left(\sum_ {s} a _ {j} \Phi \left(e _ {i}, e _ {s}\right) e _ {s}\right) = \sum_ {\tau , s} \operatorname{Tr} _ {\mathrm{A} / \mathrm{K}} \left(a _ {j} \Phi \left(e _ {i}, e _ {s}\right) a _ {\tau}\right) a _ {r} e _ {s}
\]

since \(w(ae_s) = \sum_r \mathrm{Tr}_{\Lambda/\mathbb{K}}(aa_r)a_r e_s\) for all \(a \in A\) ; thus \(u\) is the composite map \(w \circ v\) . Thus, denoting by \(v_{\kappa}\) the map \(v\) considered as a K-linear map, we have \(\det(u) = \det(v_{\kappa})\det(w)\) . Now we have \(\det(v_{\kappa}) = N_{\Lambda/\mathbb{K}}(\det(v))\) (chap.~VIII, § 12, no. 2, prop. 7), and it is clear that \(\det(v) = D_{\Phi}(e_1, \ldots, e_n)\) . On the other hand, since each of the A-modules \(Ae_i\) ( \(i = 1, \ldots, n\) ) is stable under \(w\) and the determinant of the restriction of \(w\) to \(Ae_i\) is \(\det(\mathrm{Tr}_{\Lambda/\mathbb{K}}(a_j a_r)) = D_{\Lambda/\mathbb{K}}(a_1, \ldots, a_q)\) , one has \(\det(w) = (D_{\Lambda/\mathbb{K}}(a_1, \ldots, a_q))^n\) . Formula (8) therefore reduces to \(\det(u) = \det(v_{\kappa})\det(w)\) , a formula proved above.

COROLLARY ("Transitivity formula for discriminants"). — Let K be a commutative ring, A a commutative algebra admitting a finite basis \((a_{j})_{j=1,\ldots,q}\) over K, and E an algebra over A admitting a finite basis \((e_{i})_{i=1,\ldots,n}\) . Then \((a_{j}e_{i})\) is a basis of E over K, and one has

\[
\mathrm{D} _ {\mathrm{E} / \mathrm{K}} (a _ {j} e _ {i}) = \mathrm{N} _ {\mathrm{A} / \mathrm{K}} (\mathrm{D} _ {\mathrm{E} / \mathrm{A}} (e _ {1}, \dots , e _ {n})) \cdot (\mathrm{D} _ {\mathrm{A} / \mathrm{K}} (a _ {1}, \dots , a _ {q})) ^ {n}.
\]

Indeed, if one sets \(\Phi(x,y)=\mathrm{Tr}_{\mathrm{E/A}}(xy)\) , the K-bilinear form \(\Phi'\) of Prop. 6 is \(\Phi'(x,y)=\mathrm{Tr}_{\mathrm{E/K}}(xy)\) by the transitivity formula for traces (Chap. VIII, § 12, no. 2, cor. of Prop. 7).

Exercises. — 1) Let A be an algebra of finite rank over a commutative field K, having an identity element.

\begin{enumerate}
\def\labelenumi{\alph{enumi})}
\item
  Show that if the radical of A is not zero, the bilinear form \((x, y) \to \mathrm{Tr}_{\mathbf{A}/\mathbf{K}}(xy)\) on A is degenerate.
\item
  Assume that K has characteristic 0. Show that if A is a matrix algebra \(\mathbf{M}_{n}(\mathrm{K})\) , S the canonical basis of A over K, we have \(\mathrm{D}_{\mathrm{A}/\mathrm{K}}(\mathrm{S}) \neq 0\) .
\item
  From a) and b), deduce that, for an algebra A of finite rank over a field K of characteristic 0 to be absolutely semisimple, it is necessary and sufficient that the bilinear form \((x,y)\to\mathrm{Tr}_{\mathbf{A}/\mathbf{K}}(xy)\) be nondegenerate (or, equivalently, that \(\mathrm{D}_{\mathbf{A}/\mathbf{K}}(\mathrm{S})\neq0\) for every basis S of A over K).
\end{enumerate}

Let B be a ring, A a subring of B containing the identity element of B; B is therefore an (A, A)-bimodule; we denote by \(^s\text{B}\) (resp. \(^d\text{B}\) ) the set B considered as a left (resp. right) A-module, by \(^s\text{B}^*\) (resp. \(^d\text{B}^*\) ) the right (resp. left) dual A-module of \(^s\text{B}\) (resp. \(^d\text{B}\) ). For every \(x' \in ^s\text{B}^*\) and every \(b \in \text{B}\) , \(x \to \langle xb, x' \rangle\) is an A-linear form on \(^s\text{B}\) , hence an element of \(^s\text{B}^*\) denoted by \(bx'\) ; the map \((b, x') \to bx'\) defines on \(^s\text{B}^*\) a left B-module structure (cf.~Chap.~III, 2nd ed., App. II, no. 7).

\[
(\mathrm{A}, \mathrm{A})
\]

\[
(\mathrm{A}, \mathrm{A}) -
\]

\[
\Phi : (x, y) \rightarrow \varphi (x y)
\]

\(^{s}\) B × \(^{d}\) B into A be nondegenerate, it is necessary and sufficient that \(\varphi(0)\) contain no ideal (left or right) of B other than \(\{0\}\) . We then say that \(\varphi\) is a Frobenius homomorphism from B to A.

\begin{enumerate}
\def\labelenumi{\alph{enumi})}
\setcounter{enumi}{1}
\item
  Let \(\varphi\) a Frobenius homomorphism from B to A; show that the map \(d_{\Phi}\) right-associated to \(\Phi\) is an isomorphism of the left B-module \(B_s\) onto a submodule of the left B-module \(^sB^*\) . Show that \(d_{\Phi}\) is bijective in each of the following two cases: \(1^\circ\) A is a left and right artinian ring satisfying the conditions ( \(N_s\) ) and ( \(N_d\) ) ( \(\S 1\) , exerc. 11), and \(^sB\) and \(^dB\) are free A-modules of finite length (use exerc. 11 b) of \(\S 1\) ); \(2^\circ\) A is a commutative and involutive Artinian ring (chap.~VIII, \(\S 3\) , exerc. 11), contained in the center of B, and \(^{s}\) B is an A-module of finite length (use exerc. 11 of chap.~VIII, § 3).
\item
  Conversely, show that if B \(_{s}\) and \(^{s}\) B \(^{*}\) are isomorphic, there exists a Frobenius homomorphism from B to A when one is in one of the two cases considered in b) and that \(^{s}\) B and \(^{d}\) B have the same length (use exerc. 11 b) of § 1).
\item
  With the hypotheses and notations of a), show that the right annihilator (resp. left annihilator) of a left ideal I (resp. of a right ideal r) of B is the orthogonal l⁰ (resp. r⁰) for the form Φ of the submodule I (resp. r) of sB (resp. dB).
\item
  Let \(\varphi\) a Frobenius homomorphism from B to A. Show that if A is an involutive Artinian ring (chap.~VIII, § 3, exerc. 11) then so is B, when one moreover supposes that one or the other of the conditions of \(b\) ) is satisfied (use \(b\) ) and \(d\) ).
\end{enumerate}

¶ 3) Let A be a commutative field, B an algebra of finite rank over A, having an identity element.

\begin{enumerate}
\def\labelenumi{\alph{enumi})}
\item
  For B to be a Frobenius algebra, it is necessary and sufficient that there exist a Frobenius homomorphism from B to A (exerc. 2). (Use exerc. 2 c) and e) above, and exerc. 6 b) of chap.~VIII, § 13).
\item
  Let \(\varphi\) a Frobenius homomorphism from B to A; every A-linear form on B can then be written in a unique way \(x \to \varphi(b'x)\) (resp. \(x \to \varphi(xb'')\) ) where \(b'\) and \(b''\) belong to B; for this form to be a Frobenius homomorphism, it is necessary and sufficient that \(b'\) (resp. \(b''\) ) be invertible in B.
\item
  For every \(x \in B\) , let \(x^{\sigma}\) be the unique element (cf.~b)) such that \(\varphi(xy) = \varphi(yx^{\sigma})\) for all \(y \in B\) . Show that \(x \to x^{\sigma}\) is an A-automorphism of B. One says that the Frobenius algebra B is symmetric if \(\sigma\) is an inner automorphism of B; there is then a Frobenius homomorphism from B to A for which \(\sigma\) is the identity (cf.~b)). This is equivalent to saying that the (B, B)-bimodules B and \(^{s}B^{*} = ^{d}B^{*}\) (written B*) are isomorphic (exerc. 2 c)).
\item
  Let E be a left B-module of finite length, E' its dual, E'* the dual of E' considered as a vector space over A; E'* is endowed with a structure of left B-module by setting, for \(x' \in \mathrm{E}', x'' \in \mathrm{E}*\) , \(b \in \mathrm{B}\) , \(\langle x', bx'' \rangle = \langle x'b, x'' \rangle\) (chap.~III, 2nd ed., App. II, no. 7). For every \(x \in \mathrm{E}\) , let \(f_{\mathrm{E}}(x)\) (or simply \(f(x)\) ) be the element of E'* such that \(\langle x', f(x) \rangle = \varphi(\langle x, x' \rangle)\) for all \(x' \in \mathrm{E}'\) ; show that \(f\) is a semilinear bijection for the automorphism \(\sigma\) , from the left B-module E onto the left B-module E'* (use exerc. 10 of chap.~VIII, § 4). For \(\mathrm{E} = \mathrm{B}_s\) , one has (with the notations of exerc. 2 b)) \(d_{\Phi}(x^{\sigma}) = f_{\mathrm{B}_s}(x)\) for all \(x \in \mathrm{B}\) .
\end{enumerate}

\begin{enumerate}
\def\labelenumi{\arabic{enumi})}
\setcounter{enumi}{3}
\tightlist
\item
  \begin{enumerate}
  \def\labelenumii{\alph{enumii})}
  \tightlist
  \item
    Let G be a finite group. Show that the algebra B of the group G over any commutative field A is a symmetric Frobenius algebra (exerc. 3). (Consider the map \(\varphi\) from B to A which, to every element \(x = \sum_{s \in G} \xi_s.s\) , associates \(\varphi(x) = \xi_e\) , \(e\) denoting the identity element of G).
  \end{enumerate}
\end{enumerate}

\begin{enumerate}
\def\labelenumi{\alph{enumi})}
\setcounter{enumi}{1}
\tightlist
\item
  Let E be a vector space of finite dimension n over a commutative field A, and B the exterior algebra \(\wedge\) E. Show that B is a Frobenius algebra. (If \((e_i)_{1 \leqslant i \leqslant n}\) is a basis of E, consider the map which to every element \(x\) of B associates the coefficient of \(e_1 \wedge e_2 \wedge \ldots \wedge e_n\) in the expression of \(x\) using the basis of B corresponding to \((e_i)\) ). For the Frobenius algebra B to be symmetric, it is necessary and sufficient that \(n\) be even or A have characteristic 2.
\end{enumerate}

¶ 5) a) Show that the tensor product of two Frobenius algebras (resp. Frobenius and symmetric algebras) of finite rank over a commutative field K is a Frobenius algebra (resp. Frobenius and symmetric algebra) (cf.~§ 1, no. 9, prop. 9).

\begin{enumerate}
\def\labelenumi{\alph{enumi})}
\setcounter{enumi}{1}
\tightlist
\item
  Let B be an algebra of finite rank over K, L an extension of K of finite degree over K. Show that if the algebra \(\mathrm{B}_{(\mathrm{L})} = \mathrm{B} \otimes_{\mathrm{K}} \mathrm{L}\) over L is Frobenius (resp. Frobenius and symmetric), then so is B. (Use exerc. 2 c) and 3 c) above and exerc. 2 of chap.~VIII, § 2).
\end{enumerate}

\begin{enumerate}
\def\labelenumi{\arabic{enumi})}
\setcounter{enumi}{5}
\tightlist
\item
  Show that every absolutely semisimple algebra B of finite rank over a commutative field A is a symmetric Frobenius algebra. (Reduce to the case where B is simple; use exerc. 5 b) above, as well as prop. 9 of chap.~VIII, § 12, no. 3).
\end{enumerate}

